\section{Where uncomputable content can survive}
\label{sec:states}
The fragile ruler is not the only way a computable microscopic law
can carry noncomputable asymptotic information. Two distinct homes
are a noncomputable selection among several states, and a growing
statistical readout whose expectation has a noneffective limit.
This section makes that distinction precise. All alphabets, ranges,
local tables, initial data, boundary rules and readouts called
computable have uniform algorithms. Computability of a measure means
uniform computability of its finite cylinder probabilities.
We discuss finite-alphabet classical Gibbs measures and
probabilistic cellular automata (PCA); no quantum Gibbs uniqueness
theorem is implicit.

\subsection{Effective state classes and local consensus}
\begin{lemma}[Constant observables on an effective compact class]
\label{lem:state-compact}
Let $\mathcal M$ be the probability measures on a finite-alphabet
configuration space $\Sigma^{\Z^d}$. A nonempty subset $S$ specified
by uniformly computable continuous equalities is effectively closed
in an effectively compact space. If a computable continuous local
functional $L$ is constant on $S$, that common value is computable.
In particular, a singleton $S$ is a computable measure.
\end{lemma}
\begin{proof}
For increasing finite boxes $B_n$ use the metric
$\sum_{n\geq0}2^{-n-1}\|\mu_{B_n}-\nu_{B_n}\|_{\rm TV}$.
Given an accuracy, choose $N$ making the tail small. If $B_N$ has
$M$ patterns, approximate its probability vector by a rational
simplex mesh of denominator $D$ with $M/D$ below the remaining
accuracy. Extend each finite pattern by a fixed exterior, and use
finite mixtures of these extensions as a finite net. Consistency
of marginals gives completeness and the same finite-box tests
semidecide a finite open cover. Thus the space is effectively compact.
The complement of a computable zero set is effectively open, as is
the countable union for the specified equalities.

For effectively open $V$, the inclusion $S\subset V$ is semidecidable:
search for a finite subcover of $\mathcal M$ from $V$ and the
enumerated complement of $S$. Dovetail rational intervals of width
below the requested accuracy and test $S\subset\{L\in I\}$.
If $L$ is constant, an interval containing its common value in its
interior is eventually certified. For a singleton apply this to
each cylinder probability, uniformly.
\end{proof}

This is the effective compactness mechanism behind isolated-zero
results of Galatolo--Hoyrup--Rojas
\cite[Theorem~2.9 and Corollary~3]{invariantComputability}.
The simplex mesh must have denominator depending on the number of
patterns; a fixed coarse denominator would not give the claimed net.

\begin{theorem}[Unique states and consensus]\label{thm:unique-states}
A computable finite-range PCA with a unique stationary probability
measure has a computable stationary measure. A computable bounded
finite-range classical interaction at computable finite temperature,
with a unique full DLR Gibbs measure, has a computable Gibbs measure.
More generally, any computable local expectation constant on a
nonempty effective stationary or DLR class is computable. If all
states unanimously place each specified computable local expectation
strictly on the same side of a computable threshold, the resulting
sequence of bits is computable.
\end{theorem}
\begin{proof}
For a PCA with Markov operator $P$, stationarity is given by the
countable cylinder equalities $\mu(f)=\mu(Pf)$. The finite-cone
expression for $Pf$ is computable and continuous. A stationary
measure exists by compactness of Ces\`aro averages; their stationarity
residual is at most $2/T$ on a norm-one test.
For Gibbs measures use all DLR identities
$\mu(f)=\mu(\gamma_\Lambda f)$, including tests formed by products
with exterior cylinders. Finite-range kernels are computable local
functions: their finite Boltzmann sums have strictly positive
partition functions. Compactness supplies Gibbs measures by taking
finite-box limits. These identities characterize the full DLR class;
one-site mean identities alone do not. Apply
Lemma~\ref{lem:state-compact}. Additional effectively closed
restrictions are allowed only when specified and nonemptiness is
promised. For the threshold bits dovetail the two inclusion tests
in the strict sublevel and superlevel sets. Compactness gives a
positive gap for each test separately, without a uniform gap in the
test index.
\end{proof}

The PCA uniqueness theorem is known: Marcovici--Sablik--Taati
\cite[Proposition~2.2]{noisyErgodicity} proves precisely this result.
The effective compactness and isolated-fixed-point methods of
\cite{invariantComputability} give the general framework. The
effective description of translation-invariant equilibrium-state
sets is also developed by Gayral--Sablik--Taati \cite{effectiveStateSets}.
Certified extremal local-expectation bounds of Fawzi--Fawzi--Scalet
\cite{certifiedStates} give a broader antecedent for the consensus
conclusion when the extrema coincide. The compact-cover argument
here is a direct classical specialization. Uniqueness restricted to translation-invariant measures
computes that restricted class; it must not be substituted for
uniqueness among all Gibbs measures in the following boundary result.

\begin{proposition}[What uniqueness makes effective]
\label{prop:unique-moduli}
Full Gibbs uniqueness gives a computable worst-boundary influence
modulus for each local observable, and hence an effective spatial
covariance bound for it. PCA uniqueness gives an effective uniform
\emph{Ces\`aro} convergence modulus for each local observable.
\end{proposition}
\begin{proof}
For a norm-one $f$ supported well inside $\Lambda_L$, compute the
minimum $l_L$ and maximum $u_L$ of $\gamma_{\Lambda_L}f$ over its
finite exterior collar. They enclose every infinite Gibbs expectation.
These intervals are nested by conditioning in nested boxes. If their
width failed to tend to zero, compact subsequential limits of
extremizing finite-box measures would give two Gibbs states with
different $f$ expectations. Search with rational enclosures until
$u_L-l_L<\eta$. This terminates by uniqueness. For $h$ supported
outside that box, conditioning on the exterior gives
$|\operatorname{Cov}_\mu(f,h)|\leq(u_L-l_L)\|h\|_\infty$.

For the PCA put $g_T(x)=T^{-1}\sum_{t<T}P^tf(x)$ and compute its
finite-cone minimum and maximum. Their diameter tends to zero:
otherwise subsequential limits of extremizing empirical measures
are stationary and disagree on $f$. The unique stationary expectation
lies in every such interval. Find $T$ with diameter below $\eta/2$.
For any initial law split a length-$N$ average into $T$-blocks and a
remainder. Every complete block differs from the common value by
less than $\eta/2$, and the remainder costs at most $2T/N$ for
signed norm-one $f$. Taking $N\geq4T/\eta$ gives the required
uniform modulus for all subsequent averages.
\end{proof}
The conclusion for PCA does not assert ordinary unaveraged mixing.
These algorithms can be extremely slow. An effective decay modulus
for values will also not, by itself, compute a scaling exponent.

\subsection{Noncomputable selection with computable finite preparations}
\begin{proposition}[An Ising selector]\label{prop:state-selection}
There is a computable sequence of randomized finite-volume boundary
preparations for a computable low-temperature zero-field Ising
interaction whose limiting Gibbs measure has a noncomputable local
magnetization.
\end{proposition}
\begin{proof}
Let $K$ be the halting set and put
\[
 r_s=\sum_{e\leq s}4^{-e-1}[e\text{ halts within }s\text{ steps}],
 \qquad r=\sum_{e\in K}4^{-e-1}.
\]
The limit $r\in[0,1/3]$ is noncomputable: the tail after a digit
is at most one third of that digit's weight, so a Cauchy name
recovers the halting bits successively. Choose computable
$\beta,J_x,J_y$ with
$\sinh(2\beta J_x)\sinh(2\beta J_y)>1$.
The plus and minus finite-box measures $\nu_L^\pm$ converge to
Gibbs measures $\mu^\pm$. Ferromagnetic stochastic monotonicity
gives the limits of increasing cylinder events, and finite
inclusion-exclusion gives all cylinder limits. Their magnetizations
are $\pm m$, where the Onsager--Yang formula is
\[
 m=\left[1-
  \{\sinh(2\beta J_x)\sinh(2\beta J_y)\}^{-2}\right]^{1/8}>0.
\]
For the formula and its anisotropic form see
\cite[equations~(2.1)--(2.2)]{baxterMagnetization} and
\cite{yangMagnetization}. The positive number $m$ is computable.
Choose plus boundary with the computable rational probability $r_L$
and minus otherwise. The resulting finite preparation is computable
and tends to $\mu_r=r\mu^++(1-r)\mu^-$, a DLR measure. Its local
magnetization $(2r-1)m$ computes $r$, so is noncomputable.
\end{proof}
The randomization is an explicit part of each finite preparation.
No deterministic computable infinite boundary, unique Gibbs state,
or unique selection by the Hamiltonian is claimed. The same selector
works throughout the zero-field anisotropic low-temperature family;
this is not a full interaction ball allowing a longitudinal field.

\subsection{Expectation limits and an independent-spin benchmark}
\begin{proposition}[The expectation-limit ceiling]\label{prop:state-limits}
If $\mu$ is computable and $O_s$ are uniformly computable finite
cylinder observables, any existing finite limit $\lim_s\mu(O_s)$
is in $\mathsf D_2$. A computable convergence modulus makes it
computable. A unique-state product model already allows every
limit-computable bit set as a deterministic almost-sure growing
readout limit, when the program is in the decoder.
\end{proposition}
\begin{proof}
Approximate $\mu(O_s)$ by rationals within $2^{-s}$. Their existing
limit is computable in $0'$ by tail-Cauchy search, or computable
without an oracle given a modulus. For the example take the
radius-zero PCA that independently resets every spin to
Bernoulli$(p)$, $0<p<1/2$ computable. After one step every input
law becomes the same product law, so it is the unique stationary law.
For $a(j,s)\to A(j)$ computable, take $2s+1$ distinct spins,
form their majority, and XOR it with $a(j,s)$.
The error probability is at most
$\exp[-2(2s+1)(1/2-p)^2]$, a summable sequence. The first
Borel--Cantelli lemma gives eventual agreement with the current
guess for every fixed $j$, simultaneously for countably many $j$.
Thus the readouts converge almost surely and in expectation to
$A(j)$. The bias law itself contains no such logical computation;
the decoder supplies it.
\end{proof}

The distinction between expectations and individual samples is
necessary. Under the fair computable product law,
$R_n(x)=\sum_{k=1}^n2^{-k}x_k$ has the deterministic uniform tail
bound $|R-R_n|\leq2^{-n}$, yet its sampled limit is a uniform random
real and hence noncomputable almost surely. Its expectation is
$1/2$. This random sample value is not a deterministic
noncomputable expectation. If bounded readouts converge almost
surely to a fixed constant, dominated convergence does give the
same expectation limit; an effective $L^1$ rate would compute it.

\subsection{A fault-tolerant revising computation}
\label{subsec:reliable-readout}
The next construction puts the logical revisions in a fixed
physical update rule. The reliability theorem is due to Toom and
G\'acs--Reif \cite{toom,gacsReif}; we use its arbitrary-trajectory
form as stated by Taati--Mi\k{e}kisz
\cite[Theorem~3.1]{thermalTrajectories}. The growing-readout conclusion
is an elementary corollary of that imported theorem.

\begin{theorem}[Reliable growing readouts]\label{thm:reliable-readout}
There is one fixed finite-alphabet three-dimensional stacked PCA
rule such that, for every limit-computable set $A$, a computable
cloned initialization and computable growing majority readouts
$O_{j,s}$ satisfy
\[
 O_{j,s}\longrightarrow A(j)\quad\text{almost surely},\qquad
 \mu(O_{j,s})\longrightarrow A(j)
\]
below a noise threshold. The limits persist on a relatively open
region of positive fixed-neighborhood product-update kernels.
At every computable parameter point the causal history law $\mu$
is computable. A restricted constant-normalizer thermal family gives
the same conclusion for a four-dimensional half-space DLR history
with a fixed computable bottom row.
\end{theorem}
\begin{proof}
Fix a universal one-dimensional CA simulating an ordinary finite
program for $a(j,s)\to A(j)$. Its finite initial input is computable.
Stage completion times $T_s$ and register locations $i_{j,s}$ are
computable by finite simulation until the completion marker; every
stage terminates. The same rule handles every program.
An original radius-$r$ rule can be made one-sided: compose with a
shift by $r$, then group $2r$ sites into a symbol. The blocked rule
uses only its own block and the next block. At time $T$ the original
register at $i$ is recovered at drifting coordinate $i-rT$ before
blocking. All coordinates remain computable.

Clone this one-dimensional initial tape along two redundant axes.
At each step take the full-symbol north-east-center majority in
each redundant plane (with a fixed choice for three distinct
symbols), then apply the one-dimensional rule. This is the stacking
rule of \cite[Section~3.1]{thermalTrajectories}. For any
$\delta<1/2$, its reliability theorem provides $\varepsilon_0>0$
such that joint-fault bounds $\varepsilon_0^{|F|}$ imply full-symbol
error probability at most $\delta$ at every site and time, relative
to the arbitrary noiseless trajectory. Thus revisions, rather than
only monotone outputs, are covered by the import.

For product local updates, use fresh independent coins at every
site and time. Let $R$ bound the redundant-coordinate speed. At
time $T_s$ read $m_s=2s+1$ copies of the register, separated by
$2RT_s+1$ along one redundant axis. Their finite backward cones
have disjoint update variables. Deterministic initial values do not
create a stochastic dependence. A projected bit error entails a
full-symbol error, so these independent indicators have means at
most $\delta$. Their majority obeys
\begin{equation}\label{eq:reliable-tail}
 \Pr\{O_{j,s}\ne a(j,s)\}
 \leq e^{-2m_s(1/2-\delta)^2}.
\end{equation}
For example the exponential-moment bound for centered Bernoulli
variables is $\mathbb E e^{\lambda(Y-\mathbb EY)}\leq
e^{\lambda^2/8}$; multiply over copies and optimize $\lambda$.
This proves \eqref{eq:reliable-tail} without identical distributions.
Summability and Borel--Cantelli give almost-sure eventual agreement
with each current logical guess. Countable intersection covers every
$j$. Expectation error is at most the probability in
\eqref{eq:reliable-tail} plus $|a(j,s)-A(j)|$.
Likewise $Q_s=\sum_{j\leq s}4^{-j-1}O_{j,s}$ converges almost
surely and in expectation to $\sum_{j\in A}4^{-j-1}$, by truncating
a fixed prefix and bounding the geometric tail.

A finite spacetime cylinder has a finite ancestor graph down to the
computable initial row. Summing products of computable kernel entries
over this graph computes its probability. Thus the whole causal law
is computable. Positivity and
$K(G(r)\mid r)>1-\varepsilon_0$ are finitely many strict conditions,
giving the claimed open kernel region. Conditioning first on the
past of the latest fault layer and iterating backwards bounds each
finite fault event by $\varepsilon_0^{|F|}$, even for probabilities
depending on predecessor patterns.

For the thermal version assign unit energy to a wrong update.
With $q$ symbols the kernel is
\[
 K_\beta(a\mid r)=
 \frac{e^{-\beta[a\ne G(r)]}}{1+(q-1)e^{-\beta}},\qquad
 p_\beta=\frac{(q-1)e^{-\beta}}{1+(q-1)e^{-\beta}}.
\]
For $p_\beta<\varepsilon_0$ the same readout works.
On a finite ancestor-closed directed graph the path probability is
the product of its update factors. To condition a finite resampling
region $\Lambda$, exhaust the half-space by such graphs containing
$\Lambda$ and all its immediate children. In conditional probability
ratios, precisely the factors whose support meets $\Lambda$ remain,
including factors centered at future children outside $\Lambda$.
The expression stabilizes once these children are present.
Martingale convergence proves the conditional identity for the full
exterior. Since the normalizer is independent of predecessors, this
is exactly the DLR kernel of the unit-fault interaction. This proof
uses the fixed bottom row and asserts no DLR uniqueness or stationarity.

More generally, fix for each predecessor pattern a permutation
$\pi_r$ with $\pi_r(0)=G(r)$, and give $\pi_r(k)$ the common
energy $J_k$, $J_0=0$, $J_k>0$ for $k>0$. The row normalizer
$Z=\sum_ke^{-\beta J_k}$ is still constant. The strict condition
$\sum_{k>0}e^{-\beta J_k}/Z<\varepsilon_0$ gives an open region
in this temperature/energy family. At computable points all laws
above are computable. A general energy perturbation has
predecessor-dependent normalizers, which cannot be discarded in DLR
conditioning, so arbitrary Gibbs perturbations are outside the claim.
\end{proof}

The probability-one event may depend on the kernel. Computable laws
are claimed at computable parameter points, whereas reliability is
a pointwise statement throughout the region. The effective physical
error bound controls agreement with the current guess, not the last
logical revision. The construction is a thermal history with growing
readouts; it supplies neither a noncomputable local number in a unique
equilibrium phase nor an intrinsic equilibrium critical exponent.

The history/Gibbs correspondence itself has the prior formulations of
Domany--Kinzel and Goldstein--Kuik--Lebowitz--Maes
\cite{domanyKinzel,pcaEquilibrium}, and the temperature-matched
trajectory construction in \cite{thermalTrajectories}. The direct
fixed-bottom conditional proof specifies the half-space used here.

\subsection{Scaling envelopes and the precise two-homes conclusion}
\begin{proposition}[Decay versus logarithmic scaling]\label{prop:scaling-envelope}
Effective decay of numerical correlation data does not compute its
logarithmic order. If positive computable data eventually satisfy
$c r^{-\alpha}\leq C(r)\leq C_0r^{-\alpha}$ for fixed
$0<c\leq C_0<\infty$, the exponent $\alpha$ is computable
pointwise. Uniform computation requires the envelope witnesses.
\end{proposition}
\begin{proof}
For computable rationals $\alpha_n\to\alpha\in[1,2]$, the data
$C(2^n)=2^{-n\alpha_n}\leq2^{-n}$ have an effective value-decay
bound but logarithmic order $\alpha$, possibly noncomputable.
No positive-definiteness or physical correlation realization is
asserted for this numerical example. Conversely the envelope gives
\[
 \alpha-\frac{\log_2C_0}{n}
 \leq-\frac{\log_2C(2^n)}n
 \leq\alpha-\frac{\log_2c}{n}.
\]
Hard-code an integer above the logarithmic constants and an eventual
cutoff, choose $n$ effectively large, and evaluate the positive
sample's logarithm. Fixed finite logarithmic powers add only an
effective $O(\log n/n)$ error. Uniform extraction without witnesses
fails: let $C_e(2^n)=2^{-n}$ until machine $e$ halts by time $n$,
and $2^{-2n}$ thereafter. Every sequence is eventually a pure power,
but a uniform exponent algorithm would decide halting.
\end{proof}

Here is the precise exclusion underlying the two homes. Let $S$ be
one of the nonempty effective stationary or DLR classes above,
possibly with explicitly supplied effectively closed restrictions,
and let $\mu\in S$. If a computable local observable has a
noncomputable expectation in $\mu$, its expectation is not constant
on $S$, by Lemma~\ref{lem:state-compact}. If instead $\mu$ is
computable, existing expectation limits of uniformly computable
growing cylinder observables have the ceiling and modulus condition
of Proposition~\ref{prop:state-limits}. This is a statement about
specified state classes and statistical expectations, not an
exhaustive classification of arbitrary sample-dependent observables.
State selection and fault-tolerant revision therefore coexist with
the local uniqueness obstruction.
Uniform strict contraction of a PCA instead supplies geometric
coupling decay; bounded finite-range one-dimensional finite-temperature
Gibbs chains have a uniform finite-block Doeblin contraction. In
these restricted neighborhoods the ordinary spatial correlation
length cannot diverge at an interior parameter point. Neither
argument decides the corresponding higher-dimensional equilibrium
exponent question.
